\exercisegroup

\begin{enumerate}
\def\labelenumi{\arabic{enumi}.}
\item
  A set \(\mathbf{E}\) is infinite if and only if for each mapping \(f\) of \(\mathbf{E}\) into \(\mathbf{E}\) there exists a non-empty subset \(\mathbf{S}\) of \(\mathbf{E}\) such that \(\mathbf{S} \neq \mathbf{E}\) and \(f(\mathbf{S}) \subset \mathbf{S}\) .
\item
  Show that, if \(\mathfrak{a}, \mathfrak{b}, \mathfrak{c}, \mathfrak{d}\) are four cardinals such that \(\mathfrak{a} < \mathfrak{c}\) and \(\mathfrak{b} < \mathfrak{d}\) , then \(\mathfrak{a} + \mathfrak{b} < \mathfrak{c} + \mathfrak{d}\) and \(\mathfrak{a}\mathfrak{b} < \mathfrak{c}\mathfrak{d}\) (cf.~Exercise 21 (c)).
\item
  If E is an infinite set, the set of subsets of E which are equipotent to E is equipotent to \(\mathfrak{P}(E)\) (use Proposition 3 of no. 4).
\item
  If E is an infinite set, the set of all partitions of E is equipotent to \(\mathfrak{P}(\mathbf{E})\) (associate a subset of \(E \times E\) with each partition of E).
\item
  If E is an infinite set, the set of all permutations of E is equipotent to \(\mathfrak{P}(\mathbf{E})\) . (Use Proposition 3 of no. 4 to show that, for each subset A of E whose complement does not consist of a single element, there exists a permutation f of E such that A is the set of elements of E which are invariant under f.)
\item
  Let E, F be two infinite sets such that Card (E) ≤ Card (F). Show that (i) the set of all mappings of E onto F, (ii) the set of all mappings of E into F, and (iii) the set of all mappings of subsets of E into F are all equipotent to P(F).
\item
  Let E, F be two infinite sets such that Card (E) \textless{} Card (F). Show that the set of all subsets of F which are equipotent to E and the set of all injections of E into F are both equipotent to the set \(F^{E}\) of all mappings of E into F (for each mapping f of E into F, consider the injection \(x \to (x, f(x))\) of E into \(E \times F\) ).
\item
  Show that the set of well-orderings on an infinite set E (and a fortiori the set of orderings on E) is equipotent to \(\mathfrak{P}(\mathbf{E})\) (use Exercise 5).
\item
  Let E be a non-empty well-ordered set in which every element x other than the least element of E has a predecessor (the greatest element of \(] \leftarrow, x[\) ). Show that E is isomorphic to either N or an interval \([0, n]\) of N (remark that every segment \(\neq E\) is finite by using Proposition 6 of no. 5; then use Theorem 3 of §2, no. 5).
\item
  Let \(\omega\) or \(\omega_0\) denote the ordinal Ord (N) ( \(\S 2\) , Exercise 14). The set of all integers is then a well-ordered set isomorphic to the set of all ordinals \(< \omega\) . For each integer \(n\) we denote again by \(n\) (by abuse of language) the ordinal Ord ({[}0, n{]}).
\end{enumerate}

\begin{enumerate}
\def\labelenumi{(\alph{enumi})}
\item
  Show that for each cardinal \(\mathfrak{a}\) the relation `` \(\xi\) is an ordinal and Card \((\xi) < \mathfrak{a}\) '' is collectivizing (use Zermelo's theorem). Let \(W(\mathfrak{a})\) denote the set of all ordinals \(\xi\) such that Card \((\xi) < \mathfrak{a}\) .
\item
  For each ordinal \(\alpha > 0\) define a function \(f_{\lambda}\) on the well-ordered set \(\mathrm{O}'(\alpha)\) of ordinals \(\leqslant \alpha\) by transfinite induction as follows: \(f_{\alpha}(0) = \omega_0 = \omega\) , and for each ordinal \(\xi\) such that \(0 < \xi \leqslant \alpha\) , \(f_{\alpha}(\xi)\) is the least upper bound (§ 2, Exercise 14 (d)) of the set of ordinals \(\zeta\) such that \(\operatorname{Card}(\zeta) \leqslant \operatorname{Card}(f_{\alpha}(\eta))\) for at least one ordinal \(\eta < \xi\) . Show that, if \(0 \leqslant \eta < \xi \leqslant \alpha\) , then \(\operatorname{Card}(f_{\alpha}(\eta)) < \operatorname{Card}(f_{\alpha}(\xi))\) and that, if \(\xi \leqslant \alpha \leqslant \beta\) , then \(f_{\alpha}(\xi) = f_{\beta}(\xi)\) . Put \(\omega_{\alpha} = f_{\alpha}(\alpha)\) ; \(\omega_{\alpha}\) is said to be the initial ordinal with index \(\alpha\) . We have \(\omega_{\alpha} \geqslant \alpha\) . Put \(\aleph_{\alpha} = \operatorname{Card}(\omega_{\alpha})\) ; \(\aleph_{\alpha}\) is said to be the aleph of index \(\alpha\) . In particular, \(\aleph_0 = \operatorname{Card}(\mathbf{N})\) .
\item
  Show that for each infinite cardinal \(\mathfrak{a}\) the least upper bound \(\lambda\) of the set of ordinals \(W(\mathfrak{a})\) is an initial ordinal \(\omega_{\alpha}\) , and that \(\mathfrak{a} = \aleph_{\alpha}\) (consider the least ordinal \(\mu\) such that \(\omega_{\mu} \geqslant \lambda\) ); in other words, \(\omega_{\alpha}\) is the least ordinal \(\xi\) such that \(\operatorname{Card}(\xi) = \aleph_{\alpha}\) . For each ordinal \(\alpha\) the mapping \(\xi \to \aleph_{\xi}\) , defined on \(O'(\alpha)\) , is an isomorphism of the well-ordered set \(O'(\alpha)\) onto the well-ordered set of cardinals \(\leqslant \aleph_{\alpha}\) ; in particular, \(\aleph_{\alpha+1}\) is the least cardinal \(> \aleph_{\alpha}\) . Show that, if \(\alpha\) has no predecessor, then for every strictly increasing mapping \(\xi \to \sigma_{\xi}\) of an ordinal \(\beta\) into \(\alpha\) such that \(\alpha = \sup_{\xi \in \Omega} \sigma_{\xi}\) , we have
\end{enumerate}

\[
\sum_ {\xi <   \beta} \aleph_ {\sigma_ {\xi}} = \aleph_ {\alpha}.
\]

\begin{enumerate}
\def\labelenumi{(\alph{enumi})}
\setcounter{enumi}{3}
\tightlist
\item
  Deduce from (c) that \(\omega_{\xi}\) is a normal ordinal functional symbol (§ 2, Exercise 17).
\end{enumerate}

¶ 11. (a) Show that the ordinal \(\omega\) is the least ordinal \(>0\) which has no predecessor, that \(\omega\) is indecomposable (§ 2, Exercise 16), and that for each ordinal \(\alpha > 0\) , \(\alpha\omega\) is the least indecomposable ordinal which is \(>\alpha\) (note that \(n\omega = \omega\) for each integer \(n\) ). Deduce that

\[
(\alpha + 1) \omega = \alpha \omega \quad \text { for   each } \alpha > 0.
\]

\begin{enumerate}
\def\labelenumi{(\alph{enumi})}
\setcounter{enumi}{1}
\tightlist
\item
  Deduce from (a) that an ordinal is indecomposable if and only if it is of the form \(\omega^{\beta}\) (use Exercise 18 (d) of § 2).
\end{enumerate}

\begin{enumerate}
\def\labelenumi{\arabic{enumi}.}
\setcounter{enumi}{11}
\tightlist
\item
  \begin{enumerate}
  \def\labelenumii{(\alph{enumii})}
  \tightlist
  \item
    Show that, for each ordinal \(\alpha\) and each ordinal \(\gamma > 1\) , there exist two finite sequences of ordinals \((\lambda_i)\) and \((\mu_i)\) ( \(1 \leqslant i \leqslant k\) ) such that
  \end{enumerate}
\end{enumerate}

\[
\alpha = \gamma^ {\lambda_ {1}} \mu_ {1} + \gamma^ {\lambda_ {2}} \mu_ {2} + \dots + \gamma^ {\lambda_ {k}} \mu_ {k},
\]

where \(0 < \mu_i < \gamma\) for each \(i\) , and \(\lambda_i > \lambda_{i+1}\) for \(1 \leqslant i \leqslant k - 1\) (use Exercise 18 (d) of §2 and Exercise 3 of §4). Moreover, the sequences \((\lambda_i)\) , \((\mu_i)\) are uniquely determined by these conditions. In particular, there exists a unique finite decreasing sequence \((\beta_j)_{1 \leqslant j \leqslant m}\) such that

\[
\alpha = \omega^ {\beta_ {1}} + \omega^ {\beta_ {2}} + \dots + \omega^ {\beta_ {m}}.
\]

Let \(\varphi(\alpha)\) denote the greatest ordinal \(\omega^{\beta_1}\) in this decomposition.

\begin{enumerate}
\def\labelenumi{(\alph{enumi})}
\setcounter{enumi}{1}
\tightlist
\item
  For each integer n let \(f(n) \leqslant n!\) be the greatest number of elements in the set of ordinals of the form \(\alpha_{\sigma(1)} + \alpha_{\sigma(2)} + \cdots + \alpha_{\sigma(n)}\) , where \((\alpha_i)_{1 \leqslant i \leqslant n}\) is an arbitrary sequence of n ordinals, and \(\sigma\) runs through the set of permutations of the interval {[}1, n{]}. Show that
\end{enumerate}

\[
f (n) = \sup _ {1 \leqslant k \leqslant n - 1} (k. 2 ^ {k - 1} + 1) f (n - k).\tag{1}
\]

(Consider first the case where all the \(\varphi(\alpha_i)\) are equal and show that the largest possible number of distinct ordinals of the desired form is equal to \(n\) , by using Exercise 16 (a) of §2. Then use induction on the number of ordinals \(\alpha_i\) for which \(\varphi(\alpha_i)\) takes the least possible value among the set of ordinals \(\varphi(\alpha_j)\) ( \(1 \leqslant j \leqslant n\) .) Deduce from (1) that for \(n \geqslant 20\) we have \(f(n) = 81f(n - 5)\) .

\begin{enumerate}
\def\labelenumi{(\alph{enumi})}
\setcounter{enumi}{2}
\tightlist
\item
  Show that the \(n!\) ordinals \((\omega + \sigma(1)) (\omega + \sigma(2)) \ldots (\omega + \sigma(n))\) , where \(\sigma\) runs through the set of all permutations of \([1, n]\) , are all distinct.
\end{enumerate}

¶ 13. (a) Let \(w(\xi)\) be an ordinal functional symbol (§2, Exercise 17), defined for \(\xi \geqslant \alpha_0\) and such that the relation \(\alpha_0 \leqslant \xi < \xi'\) implies \(w(\xi) < w(\xi')\) . Show that, if \(\xi \geqslant \alpha_0\) , then \(w(\xi + \eta) \geqslant w(\xi) + \eta\) for every ordinal \(\eta\) (argue by contradiction). Deduce that there exists \(\alpha\) such that \(w(\xi) \geqslant \xi\) for all \(\xi \geqslant \alpha\) (take \(\alpha\) to be the least indecomposable ordinal \(\geqslant \alpha_0\) ; cf.~Exercise 11 (a)).

\begin{enumerate}
\def\labelenumi{(\alph{enumi})}
\setcounter{enumi}{1}
\item
  Let \(f(\xi, \eta)\) be the ordinal functional symbol defined in §2, Exercise 17 (b). Suppose that the relations \(\alpha_0 \leqslant \xi \leqslant \xi'\) and \(\alpha_0 \leqslant \eta \leqslant \eta'\) imply \(g(\xi, \eta) \leqslant g(\xi', \eta')\) , so that the relations \(\alpha_0 \leqslant \xi \leqslant \xi'\) and \(1 \leqslant \eta \leqslant \eta'\) imply \(f(\xi, \eta) \leqslant f(\xi', \eta')\) (§ 2, Exercise 17 (d)). Show that for each ordinal \(\beta\) there exists at most a finite number of ordinals \(\eta\) for which the equation \(f(\xi, \eta) = \beta\) has at least one solution. (Note that if \(\xi_1\) is the least solution of \(f(\xi, \eta_1) = \beta\) and if \(\xi_2\) is the least solution of \(f(\xi, \eta_2) = \beta\) , then the relation \(\eta_1 < \eta_2\) implies \(\xi_1 > \xi_2\) .)
\item
  A critical ordinal with respect to \(f\) is any infinite ordinal \(\gamma > \alpha_0\) such that \(f(\xi, \gamma) = \gamma\) for all \(\xi\) such that \(\alpha_0 \leqslant \xi < \gamma\) . Show that a critical ordinal (with respect to \(f\) ) has no predecessor. If there exists a set \(A\) of ordinals such that \(f(\xi, \gamma) = \gamma\) for all \(\xi \in A\) , and if \(\gamma\) is the least upper bound of \(A\) , show that \(\gamma\) is a critical ordinal.
\item
  Let \(h(\xi) = f(\xi, \xi)\) (defined for \(\xi \geqslant \alpha_0\) ). Define inductively \(\alpha_1 = \alpha_0 + 2\) , \(\alpha_{n+1} = h(\alpha_n)\) for \(n \geqslant 1\) . Show that the least upper bound of the sequence \((\alpha_n)\) is a critical ordinal with respect to \(f\) .
\item
  Show that the least upper bound of every set of critical ordinals with respect to \(f\) is again a critical ordinal, and that every critical ordinal is indecomposable (note that \(f(\xi, \eta + 1) \geqslant \omega(\xi) + \eta \geqslant \xi + \eta\) for all \(\xi \geqslant \alpha_0\) ).
\end{enumerate}

¶ 14. (a) Show that if \(\alpha \geqslant 2\) and if \(\beta\) has no predecessor, then \(\alpha^{\beta}\) is an indecomposable ordinal (cf.~§ 2, Exercise 16 (a)); if \(\alpha\) is finite and if \(\beta = \omega \gamma\) , then \(\alpha^{\beta} = \omega^{\gamma}\) ; if \(\alpha\) is infinite and if \(\pi\) is the greatest indecomposable ordinal \(\leqslant \alpha\) , then \(\alpha^{\beta} = \pi^{\beta}\) (use Exercise 11).

\begin{enumerate}
\def\labelenumi{(\alph{enumi})}
\setcounter{enumi}{1}
\item
  An ordinal \(\delta\) is critical with respect to the functional symbol \(f(\xi, \eta) = \xi \eta\) if and only if, for each \(\alpha\) such that \(1 < \alpha \leqslant \delta\) , the equation \(\delta = \alpha^{\xi}\) has a solution; the unique solution \(\xi\) of this equation is then indecomposable (use Exercise 13 (e), together with Exercise 18 (d) of §2). Conversely, for each \(\alpha > 1\) and each indecomposable ordinal \(\pi\) , \(\alpha^{\pi}\) is a critical ordinal with respect to \(\xi \eta\) (use Exercise 13 (c)). Deduce that \(\delta\) is a critical ordinal with respect to \(\xi \eta\) if and only if \(\delta\) is of the form \(\omega^{\omega^{\mu}}\) (cf.~Exercise 11 (b)).
\item
  For an ordinal \(\varepsilon\) to be critical with respect to the functional symbol \(f(\xi, \eta) = \xi^{\eta}\) , i.e., such that \(\gamma^{\varepsilon} = \varepsilon\) for each \(\gamma\) satisfying \(2 \leqslant \gamma \leqslant \varepsilon\) , it is sufficient that \(2^{\varepsilon} = \varepsilon\) . Show that the least critical ordinal \(\varepsilon_0\) with respect to \(\xi^{\eta}\) is countable (cf.~Exercise 13 (d)).
\end{enumerate}

\begin{enumerate}
\def\labelenumi{\arabic{enumi}.}
\setcounter{enumi}{14}
\tightlist
\item
  Let \(\gamma\) be an ordinal \(>1\) , and for each ordinal \(\alpha\) let \(\mathbf{L}(\alpha)\) denote the set of exponents \(\lambda_{i}\) in the expression for \(\alpha\) given in Exercise 12 (a).
\end{enumerate}

\begin{enumerate}
\def\labelenumi{(\alph{enumi})}
\item
  Show that \(\lambda_{i} \leqslant \alpha\) for each \(\lambda_{i} \in \mathbf{L}(\alpha)\) , and that \(\lambda_{i} = \alpha\) for one of these ordinals only if \(\alpha = 0\) or if \(\alpha\) is a critical ordinal with respect to \(\xi^{\eta}\) (Exercise 14 (c)).
\item
  Define \(\mathbf{L}_{n}(\alpha)\) by induction on n as follows: \(\mathbf{L}_{1}(\alpha)=\mathbf{L}(\alpha)\) , and \(\mathbf{L}_{n}(\alpha)\) is the union of the sets \(\mathbf{L}(\beta)\) as \(\beta\) runs through \(\mathbf{L}_{n-1}(\alpha)\) . Show that there exists an integer \(n_{0}\) such that \(\mathbf{L}_{n+1}(\alpha)=\mathbf{L}_{n}(\alpha)\) whenever \(n\geqslant n_{0}\) , and that the elements of \(\mathbf{L}_{n}(\alpha)\) are then either 0 or critical ordinals with respect to \(\xi^{\eta}\) . (Argue by contradiction: for each n, consider the set \(\bar{\mathbf{M}}_{n}(\alpha)\) of elements \(\beta\in\mathbf{L}_{n}(\alpha)\) such that \(\beta\notin\mathbf{L}(\beta)\) , and assume that \(\mathbf{M}_{n}(\alpha)\) is not empty for any n; use (a) to obtain a contradiction.)
\end{enumerate}

\begin{enumerate}
\def\labelenumi{\arabic{enumi}.}
\setcounter{enumi}{15}
\tightlist
\item
  Every totally ordered set has a well-ordered cofinal subset (§2, Exercise 2). The least of the ordinals Ord (M) of the well-ordered cofinal subsets M of E is called the final character of E.
\end{enumerate}

\begin{enumerate}
\def\labelenumi{(\alph{enumi})}
\item
  An ordinal \(\xi\) is said to be regular if it is equal to its final character, and singular otherwise. Show that every infinite regular ordinal is an initial ordinal \(\omega_{\alpha}\) (Exercise 10). Conversely, every initial ordinal \(\omega_{\alpha}\) , whose index \(\alpha\) is either 0 or has a predecessor, is a regular ordinal. An initial ordinal \(\omega_{\alpha}\) whose index \(\alpha\) has no predecessor is singular if \(0 < \alpha < \omega_{\alpha}\) ; in particular, \(\omega_{\omega}\) is the least infinite singular initial ordinal.
\item
  An initial ordinal \(\omega_{\alpha}\) is said to be inaccessible if it is regular and its index \(\alpha\) has no predecessor. Show that, if \(\alpha = 0\) , then \(\omega_{\alpha} = \alpha\) ; in other words, \(\alpha\) is a critical ordinal with respect to the normal functional symbol \(\omega_{\eta}\) (Exercise 10 (d) and 13 (c)). Let x be the least critical ordinal with respect to this functional symbol. Show that \(\omega_{x}\) is singular, with final character \(\omega\) (cf.~Exercise 13 (d)). In other words, there exists no inaccessible ordinal \(\omega_{\alpha}\) such that \(0 < \alpha \leqslant x (*)\) .
\item
  Show that there exists only one regular ordinal which is cofinal in a given totally ordered set E; this ordinal is equal to the final character of E, and if E is not empty and has no greatest element, it is an initial ordinal. If \(\omega_{\overline{\alpha}}\) is the final character of \(\omega_{\alpha}\) , then \(\overline{\alpha} \leqslant \alpha\) ; and \(\omega_{\alpha}\) is regular if and only if \(\alpha = \overline{\alpha}\) .
\item
  Let \(\omega_{\alpha}\) be a regular ordinal and let I be a well-ordered set such that \(\operatorname{Ord}(I) < \omega_{\alpha}\) . Show that, for each family \((\xi_{i})_{i \in I}\) of ordinals such that \(\xi_{i} < \omega_{\alpha}\) for all \(i \in I\) , we have \(\sum_{i \in I} \xi_{i} < \omega_{\alpha}\) .
\end{enumerate}

\begin{enumerate}
\def\labelenumi{\arabic{enumi}.}
\setcounter{enumi}{16}
\tightlist
\item
  A cardinal \(\aleph_{\alpha}\) is said to be regular (resp. singular) if the initial ordinal \(\omega_{\alpha}\) is regular (resp. singular). For \(\aleph_{\alpha}\) to be regular it is necessary and sufficient that for every family \((a_{i})_{i\in I}\) of cardinals such that Card (I) \(< \aleph_{\alpha}\) and \(a_{i} < \aleph_{\alpha}\) for all \(i\in I\) , we have
\end{enumerate}

\[
\sum_ {i \in I} a _ {i} <   \aleph_ {\alpha}.
\]

\(\aleph_{\omega}\) is the least singular cardinal.

¶ 18. (a) For each ordinal \(\alpha\) and each cardinal \(\mathfrak{m} \neq 0\) we have \(\aleph_{\alpha+1}^{\mathfrak{m}} = \aleph_{\alpha}^{\mathfrak{m}} \cdot \aleph_{\alpha+1}\) (reduce to the case where \(\mathfrak{m} < \aleph_{\alpha+1}\) and consider the mappings of the cardinal \(\mathfrak{m}\) into the ordinal \(\omega_{\alpha+1}\) ).

\begin{enumerate}
\def\labelenumi{(\alph{enumi})}
\setcounter{enumi}{1}
\item
  Deduce from (a) that, for each ordinal \(\gamma\) such that Card \((\gamma) \leqslant m\) , we have \(\aleph_{\alpha + \gamma}^{m} = \aleph_{\alpha}^{m} \cdot \aleph_{\alpha + \gamma}^{\text{Card}(\gamma)}\) (by transfinite induction on \(\gamma\) ).
\item
  Deduce from (b) that, for each ordinal \(\alpha\) such that \(\operatorname{Card}(\alpha) \leqslant \mathfrak{m}\) , we have \(\aleph_{\alpha}^{\mathfrak{m}} = 2^{\mathfrak{m}} \cdot \aleph_{\alpha}^{\operatorname{Card}(\alpha)}\) .
\end{enumerate}

\begin{enumerate}
\def\labelenumi{\arabic{enumi}.}
\setcounter{enumi}{18}
\tightlist
\item
  \begin{enumerate}
  \def\labelenumii{(\alph{enumii})}
  \tightlist
  \item
    Let \(\alpha\) and \(\beta\) be two ordinals such that \(\alpha\) has no predecessor, and let \(\xi \to \sigma_{\xi}\) be a strictly increasing mapping of the ordinal \(\omega_{\beta}\) into the ordinal \(\alpha\) such that \(\sup \sigma_{\xi} = \alpha\) . Show that
  \end{enumerate}
\end{enumerate}

\[
\sup _ {\xi <   \omega_ {\beta}} \sigma_ {\xi} = \alpha .
\]

\[
\aleph_ {\alpha} ^ {N _ {3}} = \prod_ {\xi <   \omega_ {\beta}} \aleph_ {\sigma_ {\xi}}.
\]

(With each mapping \(f\) of the ordinal \(\omega_{\beta}\) into the ordinal \(\omega_{\alpha}\) associate an injective mapping \(\overline{f}\) of \(\omega_{\beta}\) into the set of all \(\omega_{\sigma_{\xi}}(\xi < \omega_{\beta})\) such that \(f(\zeta) \leqslant f(\zeta)\) for all \(\zeta < \omega_{\beta}\) . Calculate the cardinal of the set of mappings \(f\) associated with the same \(\overline{f}\) and observe that

\[
\mathfrak {m} = \prod_ {\xi <   \omega_ {\beta}} \aleph_ {\sigma_ {\xi}} \geqslant 2 ^ {\operatorname{Card} (\omega_ {\beta})}
\]

and \(\mathfrak{m} \geqslant \aleph_{\alpha}\) (cf.~§ 3, Exercise 3).)

\begin{enumerate}
\def\labelenumi{(\alph{enumi})}
\setcounter{enumi}{1}
\item
  Let \(\overline{\alpha}\) be the ordinal such that \(\omega_{\overline{\alpha}}\) is the final character of \(\omega_{\alpha}\) . Show that \(\aleph_{\alpha}^{\aleph_{\overline{\alpha}}} > \aleph_{\alpha}\) and that if there exists \(n\) such that \(\aleph_{\alpha} = n^{\aleph_{\gamma}}\) , then \(\gamma < \overline{\alpha}\) (use (a) and Exercise 3 of § 3).
\item
  Show that, if \(\lambda < \bar{\alpha}\) , then
\end{enumerate}

\[
\aleph_ {\alpha} ^ {\aleph_ {\lambda}} = \sum_ {\xi <   \alpha} \aleph_ {\xi} ^ {\aleph_ {\lambda}}
\]

(argue as in Exercise 18 (a)).

\begin{enumerate}
\def\labelenumi{\arabic{enumi}.}
\setcounter{enumi}{19}
\tightlist
\item
  \begin{enumerate}
  \def\labelenumii{(\alph{enumii})}
  \tightlist
  \item
    For a cardinal \(\mathfrak{a}\) to be regular (Exercise 17) it is necessary that for every cardinal \(\mathfrak{b} \neq 0\) we should have
  \end{enumerate}
\end{enumerate}

\[
a ^ {b} = a. \sum_ {m <   a} m ^ {b}.
\]

(Use Exercise 19 and consider separately the cases (i) \(\mathfrak{b}\) is finite, (ii) \(\aleph_0 \leqslant \mathfrak{b} < \mathfrak{a}\) , (iii) \(\mathfrak{b} \geqslant \mathfrak{a}\) ; also use Exercise 3 of § 3.) The generalized continuum hypothesis implies that the above condition is also sufficient.

\begin{enumerate}
\def\labelenumi{(\alph{enumi})}
\setcounter{enumi}{1}
\item
  Show that, if a cardinal \(\mathfrak{a}\) is such that \(\mathfrak{a}^{\mathfrak{m}} = \mathfrak{a}\) for every cardinal \(\mathfrak{m}\) such that \(0 < \mathfrak{m} < \mathfrak{a}\) , then \(\mathfrak{a}\) is regular (use Exercise 3 of § 3).
\item
  Show that the proposition ``for every regular cardinal a and every cardinal m such that 0 \textless{} m \textless{} a, we have \(a = a^{m}\) '' is equivalent to the generalized continuum hypothesis (use (a)).
\end{enumerate}

\begin{enumerate}
\def\labelenumi{\arabic{enumi}.}
\setcounter{enumi}{20}
\tightlist
\item
  An infinite cardinal \(\mathfrak{a}\) is said to be dominant if, for each pair of cardinals \(\mathfrak{m} < \mathfrak{a}\) , \(\mathfrak{n} < \mathfrak{a}\) , we have \(\mathfrak{m}^{\mathfrak{n}} < \mathfrak{a}\) .
\end{enumerate}

\begin{enumerate}
\def\labelenumi{(\alph{enumi})}
\item
  For \(\mathfrak{a}\) to be dominant it is sufficient that \(2^{\mathfrak{m}} < \mathfrak{a}\) for every cardinal \(\mathfrak{m} < \mathfrak{a}\) .
\item
  Define inductively a sequence \((\mathfrak{a}_n)\) of cardinals as follows: \(\mathfrak{a}_0 = \aleph_0\) , \(\mathfrak{a}_{n+1} = 2^{\mathfrak{a}_n}\) . Show that the sum \(\mathfrak{b}\) of the sequence \((\mathfrak{a}_n)\) is a dominant cardinal. \(\aleph_0\) and \(\mathfrak{b}\) are the two smallest dominant cardinals.
\item
  Show that \(\mathfrak{b}^{\aleph_0} = \aleph_0^{\mathfrak{b}} = 2^{\mathfrak{b}}\) (note that \(2^{\mathfrak{b}} \leqslant \mathfrak{b}^{\aleph_0}\) ). Deduce that \(\mathfrak{b}^{\aleph_0} = (2^{\mathfrak{b}})^{\mathfrak{b}}\) , although \(\mathfrak{b} < 2^{\mathfrak{b}}\) and \(\aleph_0 < \mathfrak{b}\) .
\end{enumerate}

\begin{enumerate}
\def\labelenumi{\arabic{enumi}.}
\setcounter{enumi}{21}
\tightlist
\item
  A cardinal \(\aleph_{\alpha}\) is said to be inaccessible if the ordinal \(\omega_{\alpha}\) is inaccessible (Exercise 16 (b)). We have then \(\omega_{\alpha} = \alpha\) if \(\omega_{\alpha} \neq \omega_{0}\) . A cardinal \(a\) is said to be strongly inaccessible if it is inaccessible and dominant (Exercise 21).
\end{enumerate}

\begin{enumerate}
\def\labelenumi{(\alph{enumi})}
\item
  The generalized continuum hypothesis implies that every inaccessible cardinal is strongly inaccessible.
\item
  For a cardinal \(\mathfrak{a} \geqslant 3\) to be strongly inaccessible it is necessary and sufficient that, for each family \((\mathfrak{a}_i)_{i \in I}\) of cardinals such that
\end{enumerate}

\[
\text { Card } (I) <   a \quad \text { and } \quad a _ {i} <   a
\]

for all \(\iota \in I\) , we should have \(\prod_{\iota \in I} a_{\iota} < a\) .

\begin{enumerate}
\def\labelenumi{(\alph{enumi})}
\setcounter{enumi}{2}
\tightlist
\item
  For an infinite cardinal \(\mathfrak{a}\) to be strongly inaccessible it is necessary and sufficient that it should be dominant (Exercise 21) and that it should satisfy one of the following two conditions: (i) \(\mathfrak{a}^{\mathfrak{b}} = \mathfrak{a}\) for every cardinal \(\mathfrak{b}\) such that \(0 < \mathfrak{b} < \mathfrak{a}\) ; (ii) \(\mathfrak{a}^{\mathfrak{b}} = \mathfrak{a}.2^{\mathfrak{b}}\) for every cardinal \(\mathfrak{b} > 0\) . (Use Exercises 20 and 21.)
\end{enumerate}

\begin{enumerate}
\def\labelenumi{\arabic{enumi}.}
\setcounter{enumi}{22}
\tightlist
\item
  Let \(\alpha\) be an ordinal \(>0\) . A mapping \(f\) of the ordinal \(\alpha\) into itself is said to be divergent if for each ordinal \(\lambda_0 < \alpha\) there exists an ordinal \(\mu_0 < \alpha\) such that the relation \(\mu_0 \leqslant \xi < \alpha\) implies \(\lambda_0 \leqslant f(\xi) < \alpha\) (*).
\end{enumerate}

\begin{enumerate}
\def\labelenumi{(\alph{enumi})}
\tightlist
\item
  Let \(\varphi\) be a strictly increasing mapping of an ordinal \(\beta\) into \(\alpha\) such that
\end{enumerate}

\[
\varphi (\sup _ {\zeta <   \gamma} \zeta) = \sup _ {\zeta <   \gamma} \varphi (\zeta) \quad \text { for   all } \quad \gamma <   \beta ,
\]

and such that

\[
\sup _ {\zeta <   \beta} \varphi (\zeta) = \alpha (\dagger).
\]

Then there exists a divergent mapping \(f\) of \(\alpha\) into itself, such that \(f(\xi) < \xi\) for all \(\xi\) satisfying \(0 < \xi < \alpha\) , if and only if there exists a divergent mapping of \(\beta\) into itself of the same type.

\begin{enumerate}
\def\labelenumi{(\alph{enumi})}
\setcounter{enumi}{1}
\item
  Deduce from (a) that there exists a divergent mapping of \(\omega_{\alpha}\) into itself, such that \(f(\xi) < \xi\) for all \(\xi\) satisfying \(0 < \xi < \alpha\) , if and only if the final character of \(\omega_{\alpha}\) is \(\omega_0\) . (If \(\omega_{\alpha}\) is a regular ordinal \(> \omega_0\) , define inductively a strictly increasing sequence \((\eta_n)\) as follows: \(\eta_1 = 1\) , and \(\eta_{n+1}\) is the least ordinal \(\zeta\) such that \(f(\xi) > \eta_n\) for all \(\xi \geqslant \zeta\) .)
\item
  Let \(\omega_{\overline{\alpha}}\) be the final character of \(\omega_{\alpha}\) (Exercise 16). Show that, if \(\overline{\alpha} > 0\) and if \(f\) is a mapping of \(\omega_{\alpha}\) into itself such that \(f(\xi) < \zeta\) for all \(\xi\) such that \(0 < \xi < \omega_{\alpha}\) , then there exists an ordinal \(\lambda_0\) such that the set of solutions of the equation \(f(\xi) = \lambda_0\) has a cardinal \(\geqslant \aleph_{\overline{\alpha}}\) .
\end{enumerate}

\begin{enumerate}
\def\labelenumi{\arabic{enumi}.}
\setcounter{enumi}{23}
\tightlist
\item
  Let \(\mathfrak{F}\) be a set of subsets of a set \(E\) such that for each \(A \in \mathfrak{F}\) we have \(\operatorname{Card}(A) = \operatorname{Card}(\mathfrak{F}) = a \geqslant \aleph_0\) . Show that \(E\) has a subset \(P\) such that \(\operatorname{Card}(P) = a\) and such that no set of \(\mathfrak{F}\) is contained in \(P\) . (If \(a = \aleph_{\alpha}\) , define by transfinite induction two injective mappings \(\xi \to f(\xi)\) , \(\xi \to g(\xi)\) of \(\omega_{\alpha}\) into \(E\) such that the sets \(P = f(\omega_{\alpha})\) and \(Q = g(\omega_{\alpha})\) do not intersect and such that each of them meets every subset \(A \in \mathfrak{F}\) .)
\end{enumerate}

\begin{enumerate}
\def\labelenumi{(\alph{enumi})}
\setcounter{enumi}{1}
\tightlist
\item
  Suppose, moreover, that for each subset \(\mathfrak{G}\) of \(\mathfrak{F}\) such that \(\operatorname{Card}(\mathfrak{G}) < \mathfrak{a}\) , the complement in \(\mathbf{E}\) of the union of the sets \(\mathbf{A} \in \mathfrak{G}\) has cardinal \(\geqslant \mathfrak{a}\) . Show that \(\mathbf{E}\) then has a subset \(\mathbf{P}\) such that \(\operatorname{Card}(\mathbf{P}) = \mathfrak{a}\) and such that, for each \(\mathbf{A} \in \mathfrak{G}\) , \(\operatorname{Card}(\mathbf{P} \cap \mathbf{A}) < \mathfrak{a}\) (similar method).
\end{enumerate}

\begin{enumerate}
\def\labelenumi{\arabic{enumi}.}
\setcounter{enumi}{24}
\tightlist
\item
  \begin{enumerate}
  \def\labelenumii{(\alph{enumii})}
  \tightlist
  \item
    Let \(\mathfrak{F}\) be a covering of an infinite set E. The degree of disjointness of \(\mathfrak{F}\) is the least cardinal \(c\) such that \(c\) is strictly greater than the cardinals Card \((\mathbf{X} \cap \mathbf{Y})\) for each pair of distinct sets \(\mathbf{X}, \mathbf{Y} \in \mathfrak{F}\) . If Card \((\mathbf{E}) = a\) and Card \((\mathfrak{F}) = b\) , show that \(b \leqslant a^c\) (note that a subset of E of cardinal \(c\) is contained in at most one set of \(\mathfrak{F}\) ).
  \end{enumerate}
\end{enumerate}

\begin{enumerate}
\def\labelenumi{(\alph{enumi})}
\setcounter{enumi}{1}
\item
  Let \(\omega_{\alpha}\) be an initial ordinal and let F be a set such that \(2 \leqslant \mathfrak{p} = \operatorname{Card}(\mathbf{F}) < \aleph_{\alpha}\) . Let E be the set of mappings of segments of \(\omega_{\alpha}\) , other than \(\omega_{\alpha}\) itself, into F. Then we have \(\operatorname{Card}(\mathbf{E}) \leqslant \mathfrak{p}^{\aleph_{\alpha}}\) . For each mapping \(f\) of \(\omega_{\alpha}\) into F, let \(\mathbf{K}_f\) be the subset of E consisting of the restrictions of \(f\) to the segments of \(\omega_{\alpha}\) (other than \(\omega_{\alpha}\) itself). Show that the set \(\mathfrak{F}\) of subsets \(\mathbf{K}_f\) is a covering of E such that \(\operatorname{Card}(\mathfrak{F}) = \mathfrak{p}^{\aleph_{\alpha}}\) , and that its degree of disjointness is equal to \(\aleph_{\alpha}\) .
\item
  Let \(\mathbf{E}\) be an infinite set of cardinal \(\mathfrak{a}\) , and let \(\mathfrak{c}\) , \(\mathfrak{p}\) be two cardinals \(> 1\) such that \(\mathfrak{p} < \mathfrak{c}\) , \(\mathfrak{p}^{\mathfrak{m}} < \mathfrak{a}\) for all \(\mathfrak{m} < \mathfrak{c}\) , and \(\mathfrak{a} = \sum_{\mathfrak{m} < \mathfrak{c}} \mathfrak{p}^{\mathfrak{m}}\) . Deduce from (b) that there exists a covering \(\mathfrak{F}\) of \(\mathbf{E}\) consisting of sets of cardinal \(\mathfrak{c}\) , with degree of disjunction equal to \(\mathfrak{c}\) , and such that \(\operatorname{Card}(\mathfrak{F}) = \mathfrak{p}^{\mathfrak{c}}\) . In particular, if \(\mathbf{E}\) is countably infinite, there exists a covering \(\mathfrak{F}\) of \(\mathbf{E}\) by infinite sets such that \(\operatorname{Card}(\mathfrak{F}) = 2^{\aleph_0}\) and such that the intersection of any two sets of \(\mathfrak{F}\) is finite.
\end{enumerate}

\begin{enumerate}
\def\labelenumi{\arabic{enumi}.}
\setcounter{enumi}{25}
\tightlist
\item
  Let E be an infinite set and let F be a set of subsets of E such that for each \(A \in F\) we have
\end{enumerate}

\[
\operatorname{Card} (\mathrm{A}) = \operatorname{Card} (\mathfrak {F}) = \operatorname{Card} (\mathrm{E}) = \mathfrak {a} \geqslant \aleph_ {0}.
\]

Show that there exists a partition \((\mathbf{B}_{\iota})_{\iota\in\mathbf{I}}\) of E such that

\[
\operatorname{Card} (\mathbf {I}) = \operatorname{Card} \left(\mathbf {B} _ {t}\right) = \mathfrak {a}
\]

for all \(\iota\in\mathbf{I}\) , and such that \(\mathbf{A}\cap\mathbf{B}_{\iota}\neq\phi\) for all \(\mathbf{A}\in\mathfrak{F}\) and all \(\iota\in\mathbf{I}\) .

(With the notation of Exercise 24 (a), consider first a surjective mapping \(f\) of \(\omega_{\alpha}\) into \(\mathfrak{F}\) such that for each \(A \in \mathfrak{F}\) the set of all \(\xi \in \omega_{\alpha}\) such that \(f(\xi) = A\) has cardinal equal to \(\mathfrak{a}\) . Then, by transfinite induction, define a bijection \(g\) of \(\omega_{\alpha}\) onto \(E\) such that \(g(\xi) \in f(\xi)\) for every \(\xi \in \omega_{\alpha}\) .)

\begin{enumerate}
\def\labelenumi{\arabic{enumi}.}
\setcounter{enumi}{26}
\tightlist
\item
  Let L be an infinite set and let \((\mathbf{E}_{\lambda})_{\lambda \in \mathbf{L}}\) be a family of sets indexed by L. Suppose that for each integer n \textgreater{} 0 the set of \(\lambda \in L\) such that \(\operatorname{Card}(\mathbf{E}_{\lambda}) > n\) is equipotent to L. Show that there exists a subset F of the product \(E = \prod_{\lambda \in L} E_{\lambda}\) , such that \(\operatorname{Card}(F) = 2^{\operatorname{Card}(L)}\) , and such that F has the following property: for each finite sequence \((f_k)_{1 \leqslant k \leqslant n}\) of distinct elements of F there exists \(\lambda \in L\) such that the elements
\end{enumerate}

\[
f _ {k} (\lambda) \in \mathrm{E} _ {\lambda} (1 \leqslant k \leqslant n)
\]

are all distinct. (Show first that there exists a partition \((\mathbf{L}_j)_{j \in \mathbb{N}}\) of \(\mathbf{L}\) such that \(\operatorname{Card}(\mathbf{L}_j) = \operatorname{Card}(\mathbf{L})\) for all \(j\) , and such that \(\operatorname{Card}(\mathbf{E}_{\lambda}) \geqslant 2^j\) for each \(\lambda \in \mathbf{L}_j\) . Hence reduce to the case where \(\mathbf{L}\) is the sum of the countable family of sets \(\mathbf{X}^j (j \geqslant 1)\) , where \(\mathbf{X}\) is an infinite set, and \(\mathbf{E}_{\lambda} = 2^j\) for each \(\lambda \in \mathbf{X}^j\) . With each mapping \(g \in 2^{\mathbf{X}}\) of \(\mathbf{X}\) into 2, associate the element \(f \in \mathbf{E}\) such that \(f(\lambda) = (g(x_1), \ldots, g(x_j))\) whenever \(\lambda = (x_k)_{1 \leqslant k \leqslant j} \in \mathbf{X}^j\) ; show that the set \(\mathbf{F}\) of elements \(f \in \mathbf{E}\) so defined has the required property.)

\begin{enumerate}
\def\labelenumi{\arabic{enumi}.}
\setcounter{enumi}{27}
\item
  Let E be an infinite set and let \((\mathcal{X}_{i})_{1\leqslant i\leqslant m}\) be a finite partition of the set \(\mathfrak{F}_{n}(\mathbf{E})\) of subsets of E having n elements. Show that there exists an index i and an infinite subset F of E such that every subset of F with n elements belongs to \(X_{i}\) . (Proof by induction on n.~For each \(a\in E\) show that there exists an index \(j(a)\) and an infinite subset \(\mathbf{M}(a)\) of \(E-\{a\}\) such that, for every subset A of \(\mathbf{M}(a)\) with n-1 elements, \(\{a\}\cup A\) belongs to \(\mathcal{X}_{j(a)}\) . Then define a sequence \((a_{i})\) of elements of E as follows: \(a_{1}\) is an arbitrary element of E, \(a_{2}\) is an arbitrary element of \(\mathbf{M}(a_{1})\) , \(a_{3}\) is defined in terms of \(\mathbf{M}(a_{1})\) and \(a_{2}\) in the same way as \(a_{2}\) was defined in terms of E and \(a_{1}\) , and so on. Show that the set F of elements of a suitable subsequence of the sequence \((a_{i})\) satisfies the required conditions.)
\item
  \begin{enumerate}
  \def\labelenumii{(\alph{enumii})}
  \tightlist
  \item
    In an ordered set E, every finite union of Noetherian subsets (with respect to the induced ordering) is Noetherian.
  \end{enumerate}
\end{enumerate}

\begin{enumerate}
\def\labelenumi{(\alph{enumi})}
\setcounter{enumi}{1}
\item
  An ordered set E is Noetherian if and only if for each \(a \in E\) the interval \(]a, \rightarrow[\) is Noetherian.
\item
  Let E be an ordered set such that the ordered set obtained by endowing E with the opposite ordering is Noetherian. Let u be a letter and let \(T\{u\}\) be a term. Show that there exists a set U and a mapping f of E onto U such that for each \(x \in E\) we have \(f(x) = T\{f^{(x)}\}\) , where \(f^{(x)}\) denotes the mapping of \(]\leftarrow, x[\) onto \(f([)\leftarrow, x[)\) which coincides with \(f\) on this interval. Furthermore, U and \(f\) are determined uniquely by this condition.
\item
  Let E be a Noetherian ordered set such that every finite subset of E has a least upper bound in E. Show that, if E has a least element, then E is a complete lattice (§ 1, Exercise 11); and that if E has no least element, the set \(E'\) obtained by adjoining a least element to E (§ 1, no. 7, Proposition 3) is a complete lattice.
\end{enumerate}

\begin{enumerate}
\def\labelenumi{\arabic{enumi}.}
\setcounter{enumi}{29}
\item
  Let E be a lattice such that the set obtained by endowing E with the opposite ordering is Noetherian. Show that every element \(a \in E\) can be written as \(\sup (e_1, e_2, \ldots, e_n)\) , where \(e_1, \ldots, e_n\) are irreducible ( \(\S 4\) , Exercise 7; show first that there exists an irreducible element e such that \(a = \sup (e, b)\) if a is not irreducible). Generalize Exercise 7(b) of \(\S 4\) to E; also generalise Exercises 8(b) and 9(b) of \(\S 4\) .
\item
  Let A be an infinite set and let E be the set of all infinite subsets of A, ordered by inclusion. Show that E is completely ramified (§ 2, Exercise 8) but not antidirected (§ 1, Exercise 23) and that E has an antidirected cofinal subset F. (Consider first the set \(\mathfrak{D}(A)\) of countable infinite subsets of A (which is cofinal in E) and let \(Z = R_0(\mathfrak{D}(A))\) (§ 1, Exercise 23). Write Z in the form \((z_{\lambda})_{\lambda \in L}\) , where L is a well-ordered set, and take F to be a set of countable subsets \(X_n^{\lambda}\) , where \(\lambda\) runs through a suitable subset of L, \(n \in N\) , \(X_m^{\lambda} \supset X_n^{\lambda}\) whenever \(m \leqslant n\) , \(X_n^{\lambda} - X_{n+1}^{\lambda}\) is infinite for all \(n \geqslant 0\) , and
\end{enumerate}

\[
\bigcap_ {n \in \mathbf {N}} \mathrm{X} _ {n} ^ {\lambda} = \phi ;
\]

the \(X_{n}^{\lambda}\) are to be defined by transfinite induction in such a way that the images of the sets \(X_{n}^{\lambda}\) under the canonical mapping \(r: \mathfrak{D}(A) \to Z\) (§1, Exercise 23) are mutually disjoint and form a cofinal subset of Z.)

¶ * 32. Let \((\mathbf{M}_n)\) , \((\mathbf{P}_n)\) be two sequences of mutually disjoint finite sets (not all empty), indexed by the set Z of rational integers. Let \(\alpha_n = \text{Card}(\mathbf{M}_n)\) , \(\beta_n = \text{Card}(\mathbf{P}_n)\) . Suppose that there exists an integer k \textgreater{} 0 such that for each \(n \in \mathbf{Z}\) and each integer \(l \geqslant 1\) we have

\[
\begin{array}{l} \alpha_ {n} + \alpha_ {n + 1} + \dots + \alpha_ {n + l} \leqslant \beta_ {n - k} + \beta_ {n - k + 1} + \dots + \beta_ {n + l + k}, \\ \beta_ {n} + \beta_ {n + 1} + \dots + \beta_ {n + l} \leqslant \alpha_ {n - k} + \alpha_ {n - k + 1} + \dots + \alpha_ {n + l + k}. \end{array}
\]

Let \(\mathbf{M}\) be the union of the family \((\mathbf{M}_n)\) and let \(\mathbf{P}\) be the union of the family \((\mathbf{P}_n)\) . Show that there exists a bijection \(\varphi\) of \(\mathbf{M}\) onto \(\mathbf{P}\) such that

\[
\varphi \left(\mathrm{M} _ {n}\right) \subset \bigcup_ {i = n - k - 1} ^ {n + k + 1} \mathrm{P} _ {i} \quad \text { and } \quad \varphi^ {- 1} \left(\mathrm{P} _ {n}\right) \subset \bigcup_ {i = n - k - 1} ^ {n + k + 1} \mathrm{M} _ {i}
\]

for each \(n \in \mathbf{Z}\) . (Consider a total ordering on each \(\mathbf{M}_n\) (resp. \(P_n\) ) and take \(\mathbf{M}\) (resp. \(P\) ) to be the ordinal sum (§ 1, Exercise 3) of the family \((\mathbf{M}_n)_{n \in \mathbf{Z}}\) (resp. \((\mathbf{P}_n)_{n \in \mathbf{Z}}\) ). If \(n_0\) is an index such that \(\mathbf{M}_{n_0} \neq \emptyset\) , consider the isomorphisms of \(\mathbf{M}\) onto \(P\) which transform the least element \(n_0 + k\)

of \(M_{n_{0}}\) into one of the elements of \(\bigcup_{j=n_{0}-k} P_{j}\) , and show that one of these isomorphisms satisfies the required conditions. Let \(\delta\) be the least of the numbers

\[
\begin{array}{l} \beta_ {n - k} + \beta_ {n - k + 1} + \dots + \beta_ {n + l + k} - (\alpha_ {n} + \alpha_ {n + 1} + \dots + \alpha_ {n + l}), \\ \alpha_ {n - k} + \alpha_ {n - k + 1} + \dots + \alpha_ {n + l + k} - (\beta_ {n} + \beta_ {n + 1} + \dots + \beta_ {n + l}) \end{array}
\]

for all \(n \in Z\) and all \(l \geqslant 1\) . If \(n \in Z\) and \(l \geqslant 1\) are such that, for example, \(\beta_{n-k} + \beta_{n-k+1} + \cdots + \beta_{n+l+k} = \delta + \alpha_n + \alpha_{n+1} + \cdots + \alpha_{n+l}\) , we may take \(\varphi\) to be such that the least element of \(P_{n-k}\) is the image under \(\varphi\) of the least element of \(M_n\) .

